\chapter{提高专题：秩、多项式与矩阵分解}\label{ch:advanced}
\chaptermap{秩不等式}{CH与矩阵幂}{最小多项式}{分解与范数}
\begin{examtip}
\textbf{竞赛拓展：}Sylvester与Frobenius不等式、矩阵幂与多项式降次、最小多项式、Jordan块计数。
\textbf{选学提高：}LU、QR、Cholesky、SVD及矩阵范数。考研复习可先掌握CH求逆和秩估计，再按需要进入其余内容。学习时先把握结论的条件，再通过核、像与不变子空间理解证明。
\end{examtip}
\section{Sylvester与Frobenius秩不等式}
\subsection{把乘积看作限制映射}
\begin{theorem}{秩不等式工具箱}
同型矩阵满足$\abs{\rank A-\rank B}\le\rank(A+B)\le\rank A+\rank B$。若$A$有$n$列、$B$有$n$行，则
\begin{equation}\rank A+\rank B-n\le\rank(AB)\le\min(\rank A,\rank B).\label{eq:sylvester}\end{equation}
对相容的$A,B,C$，Frobenius不等式为
\begin{equation}\rank(AB)+\rank(BC)\le\rank B+\rank(ABC).\label{eq:frobenius}\end{equation}
\end{theorem}
\begin{proof}
\label{der:8-1}
和矩阵的像包含于$\im A+\im B$，故$\rank(A+B)\le\dim\im A+\dim\im B$。对$A=(A+B)-B$同样使用上界，得$\rank A-\rank B\le\rank(A+B)$；交换$A,B$得到另一侧，合并即绝对值下界。

对乘积，令$X=\im B\subseteq\F^n$，限制映射$A|_X$的核是$X\cap\ker A$，像是$\im AB$。秩--零化度定理给精确等式
\[\rank AB=\rank B-\dim(X\cap\ker A).\]
由于交空间维数介于0与$\dim\ker A=n-\rank A$之间，得到$\rank AB\le\rank B$及Sylvester下界；像包含$\im AB\subseteq\im A$再给另一个上界。

对三因子，令$Y=\im BC\subseteq X=\im B$。同一限制映射分别给
\[\rank AB=\dim X-\dim(X\cap\ker A),\qquad
\rank ABC=\dim Y-\dim(Y\cap\ker A).\]
因$Y\cap\ker A\subseteq X\cap\ker A$，后一个交空间维数之差非负。两式相减得$\rank AB-\rank ABC\le\dim X-\dim Y=\rank B-\rank BC$，移项即Frobenius不等式。所有空间尺寸已在$B$的输出空间中统一，不需要方阵假设。

\end{proof}
\begin{method}{先写中间维数}
乘积下界中的减数是相乘处的公共维数，不能随便取矩阵阶数。若$AB=0$，立即用$\im B\subseteq\ker A$。若有三因子且想比较两侧乘积，优先尝试Frobenius。下界为负时以零作有效下界。
\end{method}
\worked{Sylvester界可取等}{竞赛风格}{\star\star}{8-1}
{给三阶秩均为2的$A,B$，使$\rank AB=1$。}
{让两个像方向只保留一个重合方向。}
{\textbf{方法一（直接构造）：}取$A=\diag(1,1,0)$、$B=\diag(0,1,1)$，每个各有两个非零对角元，故秩均为2。乘积$AB=\diag(0,1,0)$秩为1，达到Sylvester下界$2+2-3=1$。
\textbf{方法二（限制映射）：}$\im B=\Span(e_2,e_3)$、$\ker A=\Span(e_3)$，交维数1，所以$\rank AB=\dim\im B-\dim(\im B\cap\ker A)=2-1=1$。几何上$B$留下的两个方向，恰有一个又被$A$消去。若改取$B=A$则$AB=A$秩2，说明仅有因子秩并不能确定乘积秩。}
{构造法说明不等式下界无法普遍提高。}
{仅知道两因子秩2，乘积秩也可能为2，不能断言恒等1。}
\worked{平方秩的下界}{竞赛风格}{\star\star}{8-2}
{证明$n$阶矩阵$\rank A^2\ge2\rank A-n$，并用$J_4(0)$验证取等。}
{把Sylvester中两个因子取相同。}
{\textbf{方法一（Sylvester）：}以两个因子都取$A$，中间维数为$n$，立即得$\rank A^2\ge\rank A+\rank A-n$。
\textbf{方法二（秩损失）：}把$A$限制到$\im A$，其像为$\im A^2$、核为$\im A\cap\ker A$，所以
\[\rank A^2=\rank A-\dim(\im A\cap\ker A)
\ge\rank A-\dim\ker A=2\rank A-n.\]
验证取等：$N=J_4(0)$把$e_2,e_3,e_4$分别送到$e_1,e_2,e_3$，故秩3；$N^2$把$e_3,e_4$送到$e_1,e_2$，其余送零，故秩2。下界$2\cdot3-4=2$恰好取到。}
{核像交的维数刻画了平方损失的秩。}
{$\rank A^2$不一定等于$\rank A$，幂等与一般矩阵要区分。}
\worked{Frobenius中的中间矩阵}{竞赛风格}{\star\star}{8-3}
{取$A=(1,0),B=I_2,C=(0,1)\T$，核验Frobenius并解释为何不能省掉$\rank B$。}
{分别计算四个秩。}
{逐项乘出$AB=(1,0)$，秩1；$BC=(0,1)\T$，秩1；$ABC=(1,0)(0,1)\T=0$，秩0；而$B=I_2$秩2。因此
\[\rank AB+\rank BC=1+1=2=\rank B+\rank ABC.\]
若把右端的$\rank B$直接删去，就会要求$2\le0$，与本例矛盾，所以该项不可省。
从核像角度，$X=\im B=\R^2$，$Y=\im BC=\Span(e_2)$，且$\ker A=\Span(e_2)$。于是$\dim X-\dim Y=1$，$\dim(X\cap\ker A)-\dim(Y\cap\ker A)=0$，故$\rank AB-\rank ABC=1$。这正对应Frobenius证明中没有损失的取等情形。}
{中间空间本身携带维数，不能在乘积比较时忽略。}
{所有矩阵可为长方阵，需先核对三个乘积均有定义。}
\worked{幂等矩阵的秩分解}{竞赛风格}{\star\star\star}{8-4}
{若$P^2=P$，证明$\rank P+\rank(I-P)=n$及$\tr P=\rank P$。}
{利用互补投影的零乘积与和为单位阵。}
{\textbf{方法一（核像分解）：}$P^2=P$给$V=\im P\oplus\ker P$。令$r=\rank P$，在两空间各选一组基，拼接后$P$的矩阵为$\diag(I_r,0)$。所以$\tr P=r$，而$I-P$在像上为零、在核上恒等，故$\rank(I-P)=n-r$，两秩之和为$n$。
\textbf{方法二（谱与秩）：}特征值满足$\lambda^2=\lambda$，故为0或1；零化多项式$t(t-1)$无重根，所以$P$可对角化。设1出现$r$次，则对角形的秩与迹都为$r$，$I-P$的秩为$n-r$。两种证明均不要求$P$对称。}
{零乘积上界与和矩阵下界配对，是常用夹逼套路。}
{该迹秩等式只在幂等结构等特定条件下成立。}
\begin{insight}{乘积损失的秩，就是中间空间被核吞掉的维数}
先经$B$后经$A$，只有落在$\im B$中的向量会进入第二步；其中属于$\ker A$的部分全部变成零。因此有精确等式
\[\rank(AB)=\rank B-\dim(\im B\cap\ker A).
\]
Sylvester下界只是把交空间维数粗略地用$\dim\ker A$代替；看清交空间，通常比背不等式更接近题目的结构。
\end{insight}
\begin{figure}[H]\centering
\begin{tikzpicture}[scale=1]
\fill[Teal!8](-1.1,0)ellipse(2.1 and 1.25);\fill[Gold!10](1.1,0)ellipse(2.1 and 1.25);
\begin{scope}\clip(-1.1,0)ellipse(2.1 and 1.25);\fill[Rust!18](1.1,0)ellipse(2.1 and 1.25);\end{scope}
\draw[Teal,thick](-1.1,0)ellipse(2.1 and 1.25);\draw[Gold,thick](1.1,0)ellipse(2.1 and 1.25);
\node[lanote]at(-2.1,0){$\im B$};\node[lanote]at(2.1,0){$\ker A$};\node[lanote]at(0,0){共同部分\\被$A$消去};
\draw[-{Stealth},Ink,thick](-2,-1.55)--(2,-1.55)node[midway,below]{$\text{进入的维数}-\text{消失的维数}=\text{输出的维数}$};
\end{tikzpicture}
\caption{集合关系示意，面积不代表维数；重叠部分正是乘积秩的损失。}\label{fig:rank-intersection}
\end{figure}
\begin{clevertrick}{要取等，直接找核与像的位置}
设$A$有$n$列，记$X=\im B,K=\ker A$。限制映射$A|_X$的核为$X\cap K$、像为$\im AB$，秩--零化度给精确等式$\rank AB=\rank B-\dim(X\cap K)$。所以$\rank AB=\rank B$恰当于$\dim(X\cap K)=0$，也就是交只有零向量。
再与Sylvester下界$\rank A+\rank B-n=\rank B-\dim K$比较，二者取等恰当于
\[\dim(X\cap K)=\dim K.\]
由于$X\cap K\subseteq K$，有限维子空间同维则相等，因此此条件等价于$X\cap K=K$，即$\ker A\subseteq\im B$。这说明$A$能消去的方向全部被$B$送了进来。例\ref{ex:8-1}中$K=\Span(e_3)$、$X=\Span(e_2,e_3)$，所以交维数1，$\rank AB=2-1=1=2+2-3$，下界取等。这里指原Sylvester表达式；若其下界为负，虽可用0替代作有效下界，取等条件须重新讨论。
\textbf{易错提示：}不能只凭两因子的秩决定交空间的位置；相同的两项秩可能产生不同的乘积秩。
\end{clevertrick}


\section{凯莱--哈密顿、幂与矩阵方程}
\subsection{标准结论：将高次降到低次}
\begin{theorem}{凯莱--哈密顿定理与求逆}
若$\chi_A(t)=t^n+c_{n-1}t^{n-1}+\cdots+c_0$，则
\begin{equation}\chi_A(A)=A^n+c_{n-1}A^{n-1}+\cdots+c_0I=0.\label{eq:CH}\end{equation}
约定$c_n=1$。若$c_0\ne0$，则$A\inv=-c_0^{-1}\sum_{j=1}^n c_j A^{j-1}$，此求和式也直接包含$n=1$。任意多项式$p$对$\chi_A$作带余除法，$p=q\chi_A+r$，便有$p(A)=r(A)$。
\end{theorem}
\begin{proof}
\label{der:8-2}
写$\adj(tI-A)=\sum_{j=0}^{n-1}B_jt^j$；每个余子式次数至多$n-1$，故这种有限展开存在。代入伴随恒等式并按$t$的次数比较：
\[B_{n-1}=I,\quad B_{j-1}-AB_j=c_jI\ (1\le j\le n-1),\quad -AB_0=c_0I.\]
从最高次往下回代，得到
\[B_0=A^{n-1}+c_{n-1}A^{n-2}+\cdots+c_1I.\]
再代入常数项$-AB_0=c_0I$，把右边移到左边，正好是$A^n+c_{n-1}A^{n-1}+\cdots+c_0I=0$。$n=1$时直接由$B_0=I,-A=c_0I$得到相同结论，毋须中间递推。

为统一求逆式，记$c_n=1$。若$c_0\ne0$，提取$A$得到
\[A\left(\sum_{j=1}^n c_jA^{j-1}\right)=-c_0I.\]
括号为$A$的多项式，与$A$可交换，所以左右乘都等于$-c_0I$；乘$-\frac{1}{c_0}$即构造出双侧逆。反向$c_0=(-1)^n\det A$，非零也正是可逆条件。

若$p=q\chi_A+r$，在矩阵多项式中，由有限和与幂的交换律，$(q\chi_A)(A)=q(A)\chi_A(A)=0$，所以$p(A)=r(A)$。这说明带余除法降次的每一步都有代数依据，不依赖$A$是否可对角化。

\end{proof}
\begin{theorem}{Sylvester矩阵方程（标准结论）}
若$A\in\C^{m\times m},B\in\C^{n\times n}$的谱不交，则$AX-XB=C$对任意$C$唯一可解。若二者已对角化，进入各自特征基后逐元素有$(\lambda_i-\mu_j)y_{ij}=d_{ij}$。一般情形分别上三角化，按列递增、行递减的顺序解三角方程，系数仍为非零$\lambda_i-\mu_j$。
\end{theorem}
\begin{proof}
\label{der:sylvester-full}
复数方阵可上三角化：取一个特征向量补成基，首列只在首位有特征值，右下块再归纳，得到$A=PTP\inv,B=QSQ\inv$，其中$T,S$上三角。设$Y=P\inv XQ,D=P\inv CQ$，原方程等价于$TY-YS=D$。

逐元素写出第$(i,j)$式：
\[(t_{ii}-s_{jj})y_{ij}=d_{ij}-\sum_{k=i+1}^m t_{ik}y_{kj}+\sum_{k=1}^{j-1}y_{ik}s_{kj}.\]
按$j=1,2,\ldots,n$逐列，并在每列按$i=m,m-1,\ldots,1$向上求。右边第一和只用当前列更下面已求出的元素，第二和只用以前各列已求出的元素。谱不交保证每个$t_{ii}-s_{jj}\ne0$，因此此顺序唯一算出全部$mn$个元素。将这些元素拼成$Y$，每条方程按构造成立，再取$X=PYQ\inv$得到存在解；任何解都必须依此顺序逐项相同，故唯一。

若$A,B$已对角化，两和均为零，简化为$y_{ij}=\frac{d_{ij}}{\lambda_i-\mu_j}$。实系数情形虽然允许在复数域求解，但共轭后的解仍满足同一方程，由唯一性知解等于自身共轭，所以最后解为实矩阵。

\end{proof}
\begin{method}{先求余式再算矩阵}
若有低次零化关系，优先模该关系降次；求幂可在根上插值，重根时还需导数条件。矩阵方程若两侧分别可对角化，允许定义域和值域用不同基。易错：$AX-XB$不能一般地写成$(A-B)X$。
\end{method}
\worked{用CH求幂和逆}{考研风格}{\star\star}{8-5}
{设$A=\begin{pmatrix}1&1\\0&2\end{pmatrix}$，对非负整数$m$求$A^m$，并求$A\inv$。}
{特征多项式为$(t-1)(t-2)$，余式至多一次。}
{\textbf{方法一（CH降次）：}$\chi_A(t)=(t-1)(t-2)=t^2-3t+2$，故$A^2=3A-2I$。求$t^m$除以$(t-1)(t-2)$的余式$u_mt+v_m$，代入$t=1,2$得$u_m+v_m=1,2u_m+v_m=2^m$，所以$u_m=2^m-1,v_m=2-2^m$。于是对$m\ge0$，
\[A^m=(2^m-1)A+(2-2^m)I=\begin{pmatrix}1&2^m-1\\0&2^m\end{pmatrix}.\]
又$A(3I-A)=2I$，显式构造给$A\inv=\frac{3I-A}{2}=\begin{pmatrix}1&-\frac{1}{2}\\0&\frac{1}{2}\end{pmatrix}$。
\textbf{方法二（对角化）：}特征向量取$(1,0)\T,(1,1)\T$，令$P=\begin{pmatrix}1&1\\0&1\end{pmatrix}$，$P\inv=\begin{pmatrix}1&-1\\0&1\end{pmatrix}$。则$A=P\diag(1,2)P\inv$，把对角元分别取$m$次幂或倒数，再相乘即得上述矩阵。}
{余式法无需计算相似矩阵，特征基法更直观。}
{本题插值在互异根上即可；重根需要更多条件。}
\worked{Fibonacci矩阵}{竞赛风格}{\star\star\star}{8-6}
{令$F=\begin{pmatrix}1&1\\1&0\end{pmatrix}$，对非负整数$m$求$F^m$，并给出两条推导路线。}
{矩阵关系$F^2=F+I$与数列递推一致。}
{令$F=\begin{pmatrix}1&1\\1&0\end{pmatrix}$，Fibonacci数列满足$f_0=0,f_1=1,f_{m+1}=f_m+f_{m-1}$。
\textbf{方法一（递推归纳）：}$m=1$时$F=\begin{pmatrix}f_2&f_1\\f_1&f_0\end{pmatrix}$。若$m$时成立，右乘$F$得到
\[F^{m+1}=\begin{pmatrix}f_{m+1}+f_m&f_{m+1}\\f_m+f_{m-1}&f_m\end{pmatrix}
=\begin{pmatrix}f_{m+2}&f_{m+1}\\f_{m+1}&f_m\end{pmatrix}.\]
故所有$m\ge1$成立，$m=0$另取$I$。
\textbf{方法二（特征根）：}$\chi_F=t^2-t-1$，根$\alpha=\frac{1+\sqrt5}{2},\beta=\frac{1-\sqrt5}{2}$互异。标量递推的通解为$c\alpha^m+d\beta^m$，由$f_0=0,f_1=1$得$c+d=0,c(\alpha-\beta)=1$，故
\[f_m=\dfrac{\alpha^m-\beta^m}{\sqrt5}.\]
此数列满足同一递推与初值，归纳唯一性证明它就是$f_m$，代入矩阵公式即得显式幂。}
{CH给整数递推最简洁，谱法给闭式与增长率。}
{$m=0$应单独取$I$，不必引入未经定义的$f_{-1}$。}
\worked{幂零扰动的全部整数幂}{竞赛风格}{\star\star}{8-7}
{若$N^2=0$，求$(I+N)^m$，$m\in\mathbb Z$。}
{负整数幂先证明逆存在。}
{当$m\ge0$，因$I$与$N$可交换且$N^2=0$，二项式展开仅余常数项和一次项，得$(I+N)^m=I+mN$。当$m=-k<0$，先算$(I+N)(I-N)=I-N^2=I$，所以$(I+N)^{-1}=I-N$；再展开$(I-N)^k=I-kN=I+mN$。因此公式对所有整数成立。
另一条验证是直接设$G(m)=I+mN$，计算
\[G(m)G(\ell)=I+(m+\ell)N+m\ell N^2=G(m+\ell),\quad G(0)=I.\]
于是$G(m)^{-1}=G(-m)$，且$G(1)=I+N$，这也从乘法规律证明了整数幂公式。}
{有限多项式避免收敛问题，群关系更概括。}
{$N^3=0$时需保留$\binom m2N^2$，不能沿用只有一次项的式子。}
\worked{逐元素解Sylvester方程}{竞赛风格}{\star\star\star}{8-8}
{解$AX-XB=C$，其中$A=\diag(1,2),B=\diag(3,4)$，$C$全为1。}
{每个未知元素独立，系数为两特征值之差。}
{设$X=(x_{ij})$。左乘对角矩阵$A=\diag(1,2)$使第$i$行乘$a_i$，右乘$B=\diag(3,4)$使第$j$列乘$b_j$，故
\[(AX-XB)_{ij}=(a_i-b_j)x_{ij}.\]
四个标量方程依次为$-2x_{11}=1,-3x_{12}=1,-x_{21}=1,-2x_{22}=1$，所以
\[X=\begin{pmatrix}-\dfrac{1}{2}&-\dfrac{1}{3}\\-1&-\dfrac{1}{2}\end{pmatrix}.\]
将四项分别乘回差值均得1，验证原式。由于四个差均非零，每个标量未知数唯一，故矩阵解唯一。若某个$a_i=b_j$，该位置方程要么矛盾，要么留下自由元，唯一性就不再有保证。}
{两侧对角时无需建立四元大方程组。}
{若某个$a_i=b_j$，相应右端非零则无解，右端为零则可能不唯一。}
\section{最小多项式与Jordan块计数}
\subsection{选学提高：更小的降次依据}
\begin{definition}{零化多项式与最小多项式}
对$n$阶矩阵$A$，在$\F[t]$中取非零多项式$p$，代入时常数项按同阶单位矩阵解释。若$p(A)=0$，称$p$为$A$的零化多项式；其中次数最低的首一多项式记$m_A$。CH保证它存在。最小多项式唯一，且整除每一个零化多项式，特别整除$\chi_A$。
\end{definition}
\begin{proof}
\label{der:8-3}
CH给出至少一个非零零化多项式。把任一零化多项式除以首项系数可变成首一；这些非负整数次数中存在最小值，选该次数的首一多项式$m_A$。非零常数不可能零化非零维空间的单位阵，所以其次数至少1。

对任意零化多项式$p$作带余除法$p=qm_A+r$，$r=0$或$\deg r<\deg m_A$。代入$A$后$r(A)=p(A)-q(A)m_A(A)=0$。若$r\ne0$，把$r$首一化就得到更低次数的零化多项式，矛盾，因此$r=0$。所以$m_A$整除$p$，特别整除$\chi_A$。若$m_A,\widetilde m_A$都是首一最低次数零化多项式，彼此整除，只能相差一个非零常数；首一约定使该常数为1，故唯一。

\end{proof}
\begin{theorem}{最小多项式与结构（标准结论）}
在分裂数域上，$A$可对角化当且仅当$m_A$为互异一次因子之积。$m_A$中$(t-\lambda)$的指数等于对应最大Jordan块长度。设$d_k=\dim\ker(A-\lambda I)^k,d_0=0$，则属于$\lambda$、大小至少$k$的块数为$d_k-d_{k-1}$；恰为$k$的块数为$2d_k-d_{k-1}-d_{k+1}$。
\end{theorem}
\begin{proof}
\label{der:8-4}
Jordan形的存在与唯一已在推导\ref{der:jordan-full}由链基完整证明。现逐块分析。令$J_s(\lambda)=\lambda I+N_s$，$N_s$把第$j$个基向量送到前一向量，首向量送零，因此$\rank N_s^k=\max(s-k,0)$、$\dim\ker N_s^k=\min(k,s)$。若块属于另一根$\mu\ne\lambda$，$J_s(\mu)-\lambda I$为对角元全$\mu-\lambda\ne0$的三角矩阵，幂仍可逆，不贡献核。对所有属于$\lambda$的块求和便得$d_k$。相邻差$\min(k,s)-\min(k-1,s)$在$s\ge k$时为1，否则为0，所以至少$k$的块数为$d_k-d_{k-1}$；再减去至少$k+1$的块数，得到恰为$k$的块数$2d_k-d_{k-1}-d_{k+1}$。

把多项式在$\lambda$附近作有限展开$p(t)=\sum_{j\ge0}b_j(t-\lambda)^j$，系数$b_j=\frac{p^{(j)}(\lambda)}{j!}$可由多项式逐项求导验证。代入块后
\[p(J_s(\lambda))=b_0I+b_1N_s+\cdots+b_{s-1}N_s^{s-1}.\]
这些矩阵的非零位置分别在主对角、第1至$s-1$超对角，互不重叠，因此和为零当且仅当$b_0=\cdots=b_{s-1}=0$，也就是$p$含$(t-\lambda)^s$。对全部块同时零化，必须对每个根取最大块长度，乘起来便是最低次首一零化多项式。

所有块长为1恰当于Jordan形已对角，也恰当于最小多项式每个根指数为1。若数域不分裂，即使无重因子也未必能在该数域对角化，所以“分裂”条件仍须保留。

\end{proof}
\begin{method}{从核增长重建块}
先求$d_1,d_2$直到稳定，差分数列应非负且不增；由每个增量确定长链数，再做二阶差分得精确块数。若$m_A$已知，可先限制最大块长。不能把任意零化多项式直接当成最小多项式。
\end{method}
\worked{相同特征多项式，不同最小多项式}{竞赛风格}{\star\star}{8-9}
{求$A=I_2$与$B=\begin{pmatrix}1&1\\0&1\end{pmatrix}$的最小多项式。}
{先试一次因子是否已零化。}
{\textbf{方法一（最低次数）：}$A=I_2$满足$A-I=0$，而非零常数多项式不可能零化单位阶矩阵，故$m_A(t)=t-1$。对$B=J_2(1)$，令$N=B-I$，则$N\ne0,N^2=0$，说明$(t-1)^2$零化$B$。若有一次首一零化式$t-c$，则$B=cI$，与其非零上三角元素矛盾，因此$m_B=(t-1)^2$。
\textbf{方法二（块）：}$A$由两个大小1的根1块组成，$B$是一个大小2的根1块，最小多项式取最大块长，分别得$t-1,(t-1)^2$。它们的特征多项式同为$(t-1)^2$，说明代数重数不能替代最大块长。}
{最小多项式比特征多项式更敏感于非对角结构。}
{最小多项式次数未必等于矩阵阶数。}
\worked{多个特征值的最小多项式}{竞赛风格}{\star\star}{8-10}
{求$A=\diag(0,J_2(1))$的特征多项式和最小多项式。}
{不同根取因子之积，相同根取最大块长。}
{分块行列式给$\chi_A(t)=t\det(tI_2-J_2(1))=t(t-1)^2$。设$p(A)=0$。根0块要求$p(0)=0$。根1块写成$I+N$、$N^2=0,N\ne0$，Taylor有限展开给$p(I+N)=p(1)I+p'(1)N$；比较对角和上三角元素得$p(1)=p'(1)=0$。因此$t$与$(t-1)^2$都整除$p$，两者互素，乘积$t(t-1)^2$整除$p$。这个三次首一多项式本身确能零化各块，所以最小多项式恰为$t(t-1)^2$。}
{分块对角的最小多项式是各块最小多项式的最小公倍式。}
{若再增加一个$[1]$块，特征多项式变次，最小多项式不变。}
\worked{零化关系能告诉什么}{考研风格}{\star\star}{8-11}
{设$N^2=0,N\ne0$，求$m_N$并判断能否对角化。}
{最小次数只能在1与2之间。}
{\textbf{方法一（最小多项式）：}$N^2=0$给$m_N\mid t^2$。若$m_N=t$则$N=0$，不符条件；常数1也不能零化矩阵，因此$m_N=t^2$。它有重因子，按对角化判据不可对角化。
\textbf{方法二（反证）：}任意特征对$Nv=\lambda v$满足$0=N^2v=\lambda^2v$，而$v\ne0$，故$\lambda=0$。若$N=PDP\inv$可对角化，则$D$对角全为0，导致$N=0$，矛盾。此法不用求最小多项式，也直接证明结论。}
{低次关系可直接限制结构，通常不需行列式计算。}
{$N=0$时最小多项式为$t$，非零条件必须保留。}
\worked{六阶幂零矩阵的块}{竞赛风格}{\star\star\star}{8-12}
{幂零$N$满足$d_1=3,d_2=5,d_3=6,d_4=6$，求Jordan形与最小多项式。}
{差分为链数，最后稳定说明最大链长。}
{置$d_0=0$，令$b_k=d_k-d_{k-1}$表示大小至少$k$的Jordan块数。逐项相减得
\[(b_1,b_2,b_3,b_4)=(3,2,1,0).\]
恰为大小$k$的块数为$b_k-b_{k+1}$，所以大小1、2、3各有一个，大小至少4没有。于是
\[N\sim J_3(0)\oplus J_2(0)\oplus J_1(0),\qquad m_N(t)=t^3.\]
核验：此块形的核维数为$1+1+1=3$；平方的核维数为$2+2+1=5$；立方为$3+2+1=6$，与全部数据相符。最小多项式指数取最长块3。}
{同时检查块数总和、大小总和与核维数，可避免差分误读。}
{$d_k$在最大块长度后稳定，不一定要算到$n$次幂。}
\section{矩阵分解与范数初步}
\subsection{选学提高：结构、计算与误差}
\begin{theorem}{四种基本分解}
非奇异$n$阶矩阵经必要行置换可写$PA=LU$，其中$P$为置换矩阵，$L$为$n$阶单位下三角、$U$为$n$阶可逆上三角；若所有顺序主子式非零，可取$P=I$。满列秩实$m\times n$矩阵有薄分解$A=QR$，其中$m\ge n$、$Q\in\R^{m\times n}$满足$Q\T Q=I_n$，$R\in\R^{n\times n}$可逆上三角，可约定其对角元全正。实对称正定矩阵有唯一对角元正的Cholesky分解$A=LL\T$。任意实$m\times n$矩阵有
\begin{equation}A=U\Sigma V\T,\qquad U\T U=I_m,\quad V\T V=I_n,\label{eq:svd}\end{equation}
其中$U$、$V$分别是$m$阶、$n$阶正交矩阵；$\Sigma$为$m\times n$矩形对角矩阵，其$\min(m,n)$个对角元按$\sigma_1\ge\cdots\ge\sigma_{\min(m,n)}\ge0$排列，称奇异值。前$r=\rank A$个严格正，其余为0。
\end{theorem}
\begin{proof}
\label{der:8-5}
\textbf{LU的构造。}对非奇异$n$阶矩阵，首列必有非零元，行置换后写$P_1A=\begin{pmatrix}a&b\T\\c&D\end{pmatrix}$，$a\ne0$。消元给$S=D-\frac{cb\T}{a}$，且$\det(P_1A)=a\det S\ne0$，所以$S$非奇异。以维数归纳，存在$P_2S=L_2U_2$。于是
\[\begin{pmatrix}1&0\\0&P_2\end{pmatrix}P_1A=
\begin{pmatrix}1&0\\\dfrac{P_2c}{a}&L_2\end{pmatrix}
\begin{pmatrix}a&b\T\\0&U_2\end{pmatrix}.\]
相乘右下块为$\frac{P_2cb\T}{a}+L_2U_2=P_2D$，其余块也吻合；左因子单位下三角，右因子上三角，完成归纳，$n=1$为起点。

若全部顺序主子式$\Delta_k\ne0$，首元$a=\Delta_1$无需换行便非零。Schur补$S$的前$k$阶顺序主子式为$\frac{\Delta_{k+1}}{a}\ne0$，因为对原矩阵前$k+1$阶主块作同一次首列消元，行列式等于$a$乘该Schur子块的行列式。故对$S$归纳时也无需换行，从而整个分解无需置换。又$A$的前$k$阶主块等于相应三角主块之积，单位下三角因子行列式为1，所以$\Delta_k=\prod_{j=1}^k u_{jj}$，得到每步主元$u_{kk}=\frac{\Delta_k}{\Delta_{k-1}}\ne0$，其中$\Delta_0=1$。

无置换且单位下三角约定下，若$LU=\widetilde L\widetilde U$，则$\widetilde L\inv L=\widetilde U U\inv$同时是单位下三角与上三角，只能为$I$，故分解唯一。

\textbf{薄QR的构造。}对满列秩$A=(a_1,\ldots,a_n)$，用Schmidt定义
\[r_{ij}=q_i\T a_j\ (i<j),\quad u_j=a_j-\sum_{i<j}r_{ij}q_i,\quad r_{jj}=\norm{u_j}>0,\quad q_j=\dfrac{u_j}{r_{jj}}.\]
于是逐列$a_j=\sum_{i\le j}r_{ij}q_i$，正是$A=QR$；$Q\T Q=I$，$R$上三角且对角全正，所以可逆。每步余量非零由原列无关保证。

\textbf{Cholesky的构造与唯一。}正定$A=\begin{pmatrix}a&b\T\\b&D\end{pmatrix}$有$a>0$，配方证明$S=D-\frac{bb\T}{a}$正定：对非零$z$，取$x_1=-\frac{b\T z}{a}$，原正定二次式等于$z\T Sz>0$。归纳设$S=GG\T$且$G$对角元正，定义
\[L=\begin{pmatrix}\sqrt a&0\\\dfrac{b}{\sqrt a}&G\end{pmatrix},\qquad
LL\T=\begin{pmatrix}a&b\T\\b&\dfrac{bb\T}{a}+GG\T\end{pmatrix}=A.\]
任何正对角下三角因子的首对角元都必须是$\sqrt a$，首列其余元必须为$\frac{b}{\sqrt a}$；剔除首列后又是$S$的同类分解，归纳得到唯一。

\textbf{SVD的构造。}$A\T A$实对称半正定，取单位正交特征基$v_i$，特征值按$\lambda_1\ge\cdots\ge\lambda_r>0$与其余0排列。令$\sigma_i=\sqrt{\lambda_i}$、$u_i=\frac{Av_i}{\sigma_i}$，则
\[u_i\T u_j=\dfrac{v_i\T A\T Av_j}{\sigma_i\sigma_j}=\dfrac{\lambda_j\delta_{ij}}{\sigma_i\sigma_j}=\delta_{ij}.\]
补足为$\R^m$的单位正交基$U$。若$\lambda_j=0$，$\norm{Av_j}^2=v_j\T A\T Av_j=0$，所以$Av_j=0$。因此$AV$的前$r$列为$\sigma_i u_i$，其余全零，即$AV=U\Sigma$。右乘$V\T$得到$A=U\Sigma V\T$，且非零奇异值数等于$\rank A=r$。

\end{proof}
\begin{definition}{矩阵范数、谱半径与条件数}
向量范数满足：$\norm x\ge0$且等号仅当$x=0$；$\norm{cx}=|c|\norm x$；$\norm{x+y}\le\norm x+\norm y$。常用$\norm x_1=\sum_i|x_i|$、$\norm x_2=(\sum_i|x_i|^2)^{\frac{1}{2}}$、$\norm x_\infty=\max_i|x_i|$。

矩阵空间上的范数也满足这三条公理；方阵范数若另满足$\norm{AB}\le\norm A\norm B$，称次乘。对$A\in\R^{m\times n}$，输入、输出空间各指定向量范数后，诱导（算子）范数定义为$\norm A=\max_{x\ne0}\frac{\norm{Ax}_{\rm out}}{\norm x_{\rm in}}$。有限维下最大值存在，等价于在输入单位球面上取最大值。相乘时中间空间必须使用同一向量范数；方阵谱半径界还要求其输入输出使用同一范数。

以下$1,2,\infty$下标表示输入输出均采用相应的向量范数；Frobenius范数则按元素定义：
\begin{equation}\begin{aligned}
\norm A_1&=\max_j\sum_i|a_{ij}|,&\norm A_\infty&=\max_i\sum_j|a_{ij}|,\\
\norm A_2&=\sigma_1,&\norm A_F&=\sqrt{\sum_{ij}|a_{ij}|^2}=\sqrt{\sum_i\sigma_i^2}.
\end{aligned}\end{equation}
零矩阵的2范数为0；一般可把最大奇异值记作$\sigma_{\max}$，其中允许为零。Frobenius范数满足三条范数公理和次乘性，但一般不是诱导范数：$n\ge2$时$\norm{I_n}_F=\sqrt n$，而任一由同一向量范数诱导的$\norm{I_n}=1$。方阵的谱半径$\rho(A)=\max_i|\lambda_i|$，在复数域列出全部特征值；非奇异方阵条件数$\kappa_2(A)=\norm A_2\norm{A\inv}_2=\frac{\sigma_{\max}}{\sigma_{\min}}$。
\end{definition}
\begin{theorem}{范数与误差的基本界}
采用上述兼容的诱导范数，有$\norm{AB}\le\norm A\norm B$；对方阵及同一输入输出范数，有$\rho(A)\le\norm A$；一般$\rho(A)$不等于$\norm A_2$。非奇异$A$的$A\vx=\vb\ne0$在右端扰动下满足
\begin{equation}\dfrac{\norm{\delta\vx}_2}{\norm\vx_2}\le\kappa_2(A)\dfrac{\norm{\delta\vb}_2}{\norm\vb_2}.\end{equation}
\end{theorem}
\begin{proof}
\label{der:8-6}
\textbf{范数公理与最大值存在性。}向量1范数和无穷范数的三角不等式由$|a+b|\le|a|+|b|$逐项求和或取最大值得到；欧氏范数由Cauchy得到$\norm{x+y}_2^2\le(\norm x_2+\norm y_2)^2$，正性与齐次性都由定义成立。矩阵Frobenius范数是把全部元素排成一个向量后的欧氏范数，因此也满足这三条公理。

对任意有限维向量范数$N$，标准基展开与Cauchy给$N(x)\le\sum_i|x_i|N(e_i)\le C\norm x_2$，其中$C=(\sum_iN(e_i)^2)^{\frac{1}{2}}$。反三角不等式于是给$|N(x)-N(y)|\le C\norm{x-y}_2$，故$N$连续。在欧氏单位球面上，$N$达到严格正的最小值$c$，所以$N(x)\ge c\norm x_2$。由此$N(x)=1$的集合闭且有界，有限维下紧致；连续函数$x\mapsto N_{\rm out}(Ax)$在其上达到最大值，证明诱导范数中的最大值存在。由这个定义逐项使用向量范数公理，可得诱导范数的正性、齐次性与三角不等式；正性中若$A\ne0$，选非零列对应的$e_j$便有正的长度比。

\textbf{先推导各范数公式。}以下$A$为实矩阵。对1范数，
\[\norm{Ax}_1\le\sum_i\sum_j|a_{ij}||x_j|\le\left(\max_j\sum_i|a_{ij}|\right)\norm x_1.\]
取最大列对应的$x=e_j$达到上界，所以诱导1范数等于最大列和。对无穷范数，$|(Ax)_i|\le\sum_j|a_{ij}|\norm x_\infty$，故上界为最大行和。取最大行$i_0$，实数域令$x_j=\sgn(a_{i_0j})$（零项任取绝对值1），便有$|(Ax)_{i_0}|=\sum_j|a_{i_0j}|$且$\norm x_\infty=1$，上界达到；复数域取相反相位同样达到。

SVD中令$x=Vy$，正交性保长度，故$\norm{Ax}_2^2=\sum_{i\le r}\sigma_i^2y_i^2\le\sigma_1^2\sum_jy_j^2$。取$y=e_1$达到上界，所以$\norm A_2=\sigma_1$；零矩阵时取0。又
\[\sum_{ij}|a_{ij}|^2=\tr(A\T A)=\tr(V\Sigma\T\Sigma V\T)=\tr(\Sigma\T\Sigma)=\sum_i\sigma_i^2,\]
中间使用$\tr(XY)=\tr(YX)$，因此得到Frobenius公式。对相容实矩阵，逐个$(i,j)$位置用Cauchy并求和，得到
\[\begin{aligned}
\norm{AB}_F^2&=\sum_{i,j}\left|\sum_k a_{ik}b_{kj}\right|^2\\
&\le\sum_{i,j}\left(\sum_k|a_{ik}|^2\right)\left(\sum_k|b_{kj}|^2\right)=\norm A_F^2\norm B_F^2.
\end{aligned}\]
开平方即Frobenius次乘性；它不要求该范数是诱导范数。非奇异方阵逆的SVD为$V\Sigma\inv U\T$，最大奇异值是$\frac{1}{\sigma_{\min}}$，得到条件数比值公式。

\textbf{次乘性与谱界。}对任意非零$x$，$\norm{ABx}\le\norm A\norm{Bx}\le\norm A\norm B\norm x$，取比值的最大值得次乘性。在复数域定义的向量范数中，对非零特征向量直接得到$|\lambda|\norm z=\norm{Az}\le\norm A\norm z$，除以$\norm z>0$便得谱界。

若原矩阵与向量范数只在实数域定义，而特征向量$z=u+iv$为复数，须说明如何使用它。定义
\[\mathcal N(z)=\max_{0\le\theta\le2\pi}\norm{\operatorname{Re}(e^{i\theta}z)}.\]
连续函数在闭区间达到最大值；若$z\ne0$，则$u,v$至少一个非零，取$\theta=0$或$\frac{\pi}{2}$说明$\mathcal N(z)>0$。实范数三角不等式逐个$\theta$成立，再取最大值得$\mathcal N(z+w)\le\mathcal N(z)+\mathcal N(w)$。对$\alpha=re^{i\phi}$，其中$r\ge0$，实范数的齐次性与角度平移给
\[\mathcal N(\alpha z)=r\max_\theta\norm{\operatorname{Re}(e^{i(\theta+\phi)}z)}=|\alpha|\mathcal N(z).\]
所以$\mathcal N$确为复向量范数，且对实$z$等于原范数。实矩阵$A$与取实部可交换，因此逐个$\theta$有
\[\norm{\operatorname{Re}(e^{i\theta}Az)}
=\norm{A\operatorname{Re}(e^{i\theta}z)}
\le\norm A\,\norm{\operatorname{Re}(e^{i\theta}z)}.\]
取最大值得$\mathcal N(Az)\le\norm A\,\mathcal N(z)$。再将$Az=\lambda z$代入并除以$\mathcal N(z)>0$，仍得$|\lambda|\le\norm A$，因而所有实数域诱导范数均满足$\rho(A)\le\norm A$。

\textbf{右端误差。}保持$A$不变，$A(x+\delta x)=b+\delta b$减去$Ax=b$得$\delta x=A\inv\delta b$，所以$\norm{\delta x}\le\norm{A\inv}\norm{\delta b}$。同时$\norm b\le\norm A\norm x$，由于$b\ne0$可除，得到
\[\dfrac{\norm{\delta x}}{\norm x}\le\norm{A\inv}\norm{\delta b}\dfrac{\norm A}{\norm b}=\kappa(A)\dfrac{\norm{\delta b}}{\norm b}.\]
取欧氏诱导范数即为正文公式。这里仅扰动右端；若连$A$一起扰动，必须另写误差方程，不能直接沿用此式。

\end{proof}
\begin{method}{分解服务于目标}
解多右端线性系统优先LU；最小二乘优先QR；正定系统用Cholesky；秩、最小范数与几何压缩用SVD。估计误差时先选范数与尺度。范数小不代表可逆，行列式小也不单独代表条件数大。
\end{method}
\worked{LU与Cholesky比较}{考研风格}{\star\star}{8-13}
{分解$A=\begin{pmatrix}2&1\\1&2\end{pmatrix}$，给LU与Cholesky。}
{消元主元为2与$\frac{3}{2}$。}
{\textbf{方法一（LU）：}对第二行减去第一行的$\frac{1}{2}$，得到$U=\begin{pmatrix}2&1\\0&\frac{3}{2}\end{pmatrix}$。消元矩阵$E=\begin{pmatrix}1&0\\-\frac{1}{2}&1\end{pmatrix}$满足$EA=U$，故$A=E\inv U=LU$，其中$L=\begin{pmatrix}1&0\\\frac{1}{2}&1\end{pmatrix}$。相乘第二行得$(1,\frac{1}{2}+\frac{3}{2})=(1,2)$，验证成立。
\textbf{方法二（Cholesky）：}设$G=\begin{pmatrix}a&0\\b&c\end{pmatrix}$，约定$a,c>0$。比较$GG\T=A$得到$a^2=2,ab=1,b^2+c^2=2$，依次解为$a=\sqrt2,b=\frac{1}{\sqrt2},c=\sqrt{\frac{3}{2}}$。正主子式$2,3$保证正定和根号合法；正对角约定使三个数逐项唯一。}
{Cholesky利用对称正定，存储与运算更简洁；LU适用范围更广。}
{Cholesky下三角的对角元取正，不要把单位下三角$L$直接当作平方根因子。}
\worked{正交列的QR}{基础例题}{\star\star}{8-14}
{求$A=\begin{pmatrix}1&1\\1&-1\end{pmatrix}$的QR分解。}
{两列已正交，只需单位化。}
{\textbf{方法一（正交列单位化）：}两列$a_1=(1,1)\T,a_2=(1,-1)\T$满足$a_1\T a_2=0$、长度均为$\sqrt2$。故
\[Q=\dfrac1{\sqrt2}\begin{pmatrix}1&1\\1&-1\end{pmatrix},\qquad R=\sqrt2 I_2.\]
计算$Q\T Q=\frac{A\T A}{2}=I$，又$QR=A$，所以是QR分解。
\textbf{方法二（Schmidt系数）：}$r_{11}=\sqrt2,q_1=\frac{a_1}{\sqrt2}$；$r_{12}=q_1\T a_2=\frac{1-1}{\sqrt2}=0$，故$u_2=a_2$，$r_{22}=\sqrt2,q_2=\frac{a_2}{\sqrt2}$。按$r_{ij}$拼出的上三角$R$与前法相同。$\det Q=-1$仅说明含反射，不影响正交性。}
{已正交时无需重复完整消元。}
{$Q$的行列式为$-1$仍是正交矩阵，不必强求为$1$。}
\worked{矩形SVD与伪逆}{竞赛风格}{\star\star}{8-15}
{求$A=\begin{pmatrix}3&0\\0&1\\0&0\end{pmatrix}$的SVD、2范数、Frobenius范数与最小二乘系数。}
{矩形对角形式已是SVD。}
{取$U=I_3,V=I_2,\Sigma=A$，则$U\Sigma V\T=A$，两侧均正交，矩形对角非负且降序，所以已经是SVD。$A\T A=\diag(9,1)$，奇异值为3、1。因此$\norm A_2=3$，$\norm A_F=\sqrt{3^2+1^2}=\sqrt{10}$。
对任意$\vb=(b_1,b_2,b_3)\T$，误差平方为
\[\norm{A\vx-\vb}^2=(3x_1-b_1)^2+(x_2-b_2)^2+b_3^2.\]
前两项同时为零给唯一最小二乘系数$(\frac{b_1}{3},b_2)\T$，残差为$(0,0,b_3)\T$，最小误差$|b_3|$。把该系数写成$A^+\vb$即得
\[A^+=\begin{pmatrix}\dfrac{1}{3}&0&0\\0&1&0\end{pmatrix}.\]
可检验$A^+A=I_2,AA^+=\diag(1,1,0)$，后者恰是到列空间的正交投影。}
{SVD把最佳逼近拆成各个独立缩放方向。}
{矩形$\Sigma$不是方阵，不能写普通逆$\Sigma\inv$。}
\worked{谱半径、范数与条件数}{竞赛风格}{\star\star\star}{8-16}
{设$M>0$，比较$N=\begin{pmatrix}0&M\\0&0\end{pmatrix}$的谱半径与2范数；另求$D_\varepsilon=\diag(1,\varepsilon)$的条件数，$0<\varepsilon\le1$。}
{特征值与奇异值度量不同现象。}
{\textbf{方法一（特征值与奇异值）：}$N$上三角且对角为0，故全部特征值为0，$\rho(N)=0$。而$N\T N=\diag(0,M^2)$的最大特征值为$M^2$，因此$\norm N_2=M$。$D_\varepsilon\T D_\varepsilon=\diag(1,\varepsilon^2)$，因$0<\varepsilon\le1$，最大、最小奇异值为$1,\varepsilon$，得$\kappa_2(D_\varepsilon)=\frac{1}{\varepsilon}$。
\textbf{方法二（直接长度比）：}对非零$(x,y)$，$\frac{\norm{N(x,y)\T}_2}{\sqrt{x^2+y^2}}=\frac{M|y|}{\sqrt{x^2+y^2}}\le M$，在$x=0$时取等。$D_\varepsilon$沿两坐标轴的缩放率为1、$\varepsilon$，所有单位向量的长度平方$x^2+\varepsilon^2y^2$夹在$\varepsilon^2$与1之间，故最大最小缩放率的比为$\frac{1}{\varepsilon}$。例如右端沿第一轴而扰动沿第二轴时，这种缩放差会放大相对解误差。}
{谱半径描述特征方向，诱导范数描述所有方向的最大瞬时缩放。}
{小特征值或行列式不能脱离其他尺度单独衡量病态；$\varepsilon I$条件数恒为1。}
\begin{insight}{Jordan链解释核为何逐级增长}
在一条长度$s$的幂零链上，一次作用把核端的首向量$v_1$送到零，$k$次作用能消去$\min(k,s)$个独立方向。把各条链的贡献相加，就得到核维数。增长量不增，是因为“长度至少$k+1$”的链集合包含于“长度至少$k$”的链集合；这不表示恰为某长度的块数必须单调。
\end{insight}
\begin{figure}[H]\centering
\begin{tikzpicture}[scale=1.05]
\foreach \xx/\txt in {0/v_3,1.2/v_2,2.4/v_1}{\node[draw=Teal,fill=Teal!6,minimum width=8mm]at(\xx,2){$\txt$};}
\foreach \xx in {0.4,1.6,2.8}{\draw[-{Stealth},Teal](\xx,2)--(\xx+.4,2);}\node at(3.6,2){$0$};
\node[draw=Gold,fill=Gold!6,minimum width=8mm]at(1.2,1){$w_2$};\node[draw=Gold,fill=Gold!6,minimum width=8mm]at(2.4,1){$w_1$};
\draw[-{Stealth},Gold](1.6,1)--(2,1);\draw[-{Stealth},Gold](2.8,1)--(3.2,1);\node at(3.6,1){$0$};
\node[draw=Ink,fill=Ink!4,minimum width=8mm]at(2.4,0){$z_1$};\draw[-{Stealth},Ink](2.8,0)--(3.2,0);\node at(3.6,0){$0$};
\node[lanote]at(1.8,-.75){三条链的长度为$3,2,1$};
\draw[laaxis](5,-.05)--(8.5,-.05);
\foreach \xx/\hh/\num/\lab in {5.5/.9/3/1,6.6/1.5/5/2,7.7/1.8/6/3}{\filldraw[fill=Teal!15,draw=Teal](\xx-.3,0)rectangle(\xx+.3,\hh);\node[above]at(\xx,\hh){$\num$};\node[font=\small,below]at(\xx,-.12){$d_{\lab}$};}
\node[lanote]at(6.6,-.75){核维数：$3,5,6$};
\end{tikzpicture}
\caption{每一条箭头表示作用一次$N$；核增长量为$3,2,1$，依次是长度至少1、2、3的链数。}\label{fig:jordan-kernel-growth}
\end{figure}
\begin{clevertrick}{看核的差分，不猜Jordan形}
给出$d_0=0,d_1=3,d_2=5,d_3=6$且此后稳定，令$b_k=d_k-d_{k-1}$，逐项得$b_1=3,b_2=2,b_3=1,b_4=0$。一条长度$s$的链对$d_k$贡献$\min(k,s)$，所以对$b_k$贡献1恰在$s\ge k$时成立；故$b_k$数的是长度至少$k$的块。恰为长度$k$的块数则是$b_k-b_{k+1}$，本例为
\[b_1-b_2=1,\qquad b_2-b_3=1,\qquad b_3-b_4=1.\]
所以块为$J_1(0),J_2(0),J_3(0)$各一个。总大小$1+2+3=6$；反算核维数为$1+1+1=3,1+2+2=5,1+2+3=6$，验证全部数据，无需枚举六阶矩阵的所有分拆。
\textbf{易错提示：}核维数必须不减，增长量必须不增；若计算得$d_1=1,d_2=3$，就与链计数矛盾，应返回检查消元。
\end{clevertrick}
\begin{insight}{SVD是“换正交方向、分别缩放、再换方向”}
$A=U\Sigma V\T$中，$V\T$先转到输入主方向，$\Sigma$按奇异值缩放，$U$再把输出主方向摆回去。对二维可逆矩阵，单位圆像是非退化椭圆，其最长、最短半轴为两个正奇异值；若矩阵奇异，图形会退化为线段或一点，不能仍按两个正半轴读取。非奇异方阵的条件数比较最大与最小缩放率，而不是面积。
\end{insight}
\begin{figure}[H]\centering
\begin{tikzpicture}[scale=1.15]
\begin{scope}
\draw[laaxis](-1.4,0)--(1.5,0);\draw[laaxis](0,-1.4)--(0,1.5);\draw[Teal,thick](0,0)circle(1);
\node[lanote]at(0,-1.85){输入单位圆};
\end{scope}
\draw[-{Stealth},Gold,thick](1.8,0)--(3.2,0)node[midway,above]{$\diag(3,1)$};
\begin{scope}[xshift=6.8cm]
\draw[laaxis](-3.4,0)--(3.5,0);\draw[laaxis](0,-1.4)--(0,1.5);\draw[Teal,thick](0,0)ellipse(3 and 1);
\draw[lavec,Ink](0,0)--(3,0)node[midway,below]{$\sigma_1=3$};\draw[lavec,Gold](0,0)--(0,1)node[above left]{$\sigma_2=1$};
\node[lanote]at(0,-1.85){输出椭圆，$\kappa_2=3$};
\end{scope}
\end{tikzpicture}
\caption{奇异值是单位球各主方向的缩放率；一般矩阵还会在两侧作正交方向变换。}\label{fig:svd-ellipse}
\end{figure}
\begin{clevertrick}{用尺度反例区分行列式和条件数}
先取$0<\varepsilon<1$，$A=\varepsilon I_2$。有$\det A=\varepsilon^2$、$A\T A=\varepsilon^2I_2$，所以两个奇异值都为$\varepsilon$，$\kappa_2(A)=\frac{\varepsilon}{\varepsilon}=1$。当$\varepsilon\to0^+$，行列式趋于零，但在每个非零$\varepsilon$下条件数都为1。
再取$M\ge1$，$B=\diag(M,\frac{1}{M})$。有$\det B=M\cdot(\frac{1}{M})=1$，$B\T B=\diag(M^2,M^{-2})$，所以最大、最小奇异值为$M,\frac{1}{M}$，
\[\kappa_2(B)=\dfrac{M}{\frac{1}{M}}=M^2\longrightarrow\infty\quad(M\to\infty).\]
前者在各方向均匀缩小，后者沿一轴拉长、另一轴压缩；面积分别记录奇异值的乘积，而条件数比较最大与最小的比值。图中$\diag(3,1)$同理有奇异值3、1，故$\kappa_2=3$，不能仅凭面积判断方向伸缩是否均衡。
\end{clevertrick}


\begin{warning}
CH可用于不可对角化矩阵。最小多项式必须首一且次数最小。Jordan存在需特征多项式分裂，实数域可能需要转到复数域。QR的$R$可逆依赖满列秩，Cholesky依赖正定。奇异值不是一般矩阵特征值的绝对值，只有正规矩阵等特殊情形才有这种关系。
\end{warning}
\section*{章末分层习题}\addcontentsline{toc}{section}{章末分层习题}
\exercise{8-1}{基础}{五阶$A,B$秩为4、3，求$\rank AB$的上下界。}
\exercise{8-2}{基础}{若$N^2=0$，计算$(I+2N)^5$。}
\exercise{8-3}{基础}{求$\diag(1,1,2)$的最小多项式。}
\exercise{8-4}{考研风格}{用CH求$\begin{pmatrix}2&1\\1&1\end{pmatrix}$的逆。}
\exercise{8-5}{考研风格}{四阶幂零矩阵$\rank N=2,\rank N^2=1,\rank N^3=0$，求Jordan形。}
\exercise{8-6}{考研风格}{求$\begin{pmatrix}1&-2\\3&4\end{pmatrix}$的1范数、无穷范数与Frobenius范数。}
\exercise{8-7}{竞赛拓展}{证明$A$的最小多项式整除任意零化多项式，并据此说明唯一性。}
\exercise{8-8}{竞赛拓展}{设$A$可逆，证明$\kappa_2(A)\ge1$，并刻画等号。}
\exercise{8-9}{竞赛拓展}{若$\norm E<1$（诱导范数），证明$I-E$可逆且$\norm{(I-E)\inv}\le\frac{1}{1-\norm E}$。}
\section*{章末方法速查}\addcontentsline{toc}{section}{章末方法速查}
\begin{quicktable}
核像夹逼&乘积与零乘积&秩的上下界&中间维数准确\\
Frobenius&三个相容矩阵&左右乘积同时出现&注意不等号方向\\
CH降次&已知特征多项式&矩阵幂与逆&标量常数写成$I$\\
最小多项式&需最低次数关系&重复谱、幂零&零化不等于最小\\
核差分&分裂数域&Jordan块大小&差分应非增\\
分解与范数&数值或几何目标&解方程、投影、误差&核对条件和尺度\\
\end{quicktable}
